

%%%%%%%%%%%%%%%%%%%%%%%%%
\section{Diagrammatics for Soergel Bimodules} \label{sec-diagsbim}
%%%%%%%%%%%%%%%%%%%%%%%%%

In this chapter we define a diagrammatic category $\DC$ by generators and relations. We provide a functor $\FC$ from this category to Bott-Samelson bimodules. The main result of this paper
is that $\FC$ is an equivalence, and the proof will be given in later chapters.

Fix a realization $\hg$ of $(W,S)$ over $\Bbbk$, as in section \ref{subsec-realization}, with $S$ finite (though see Remark \ref{SCnotfinite}). As before, we let $R = \Bbbk[\hg]$ denote
the coordinate ring of the reflection representation, $\a_s$ the simple root associated to $s \in S$, $\pa_s$ the Demazure operator. We let $Q$ denote the field of fractions of $R$.

For reasons of simplicity, we will assume the realization is balanced in the first pass. Later, in section \ref{nonbalanced}, we treat the unbalanced case. For reasons to become clear in
that section, we must always assume our realization is even-balanced.

%===========
\subsection{Generators and relations}
%===========

\begin{defn} A \emph{Soergel graph} is a type of graph with boundary (see Definition \ref{defn: graph}). The only decorations are boxes labelled by homogeneous $f \in R$. The vertices in this graph are of 3 types (see Figure \ref{thegenerators}):

\begin{itemize} \item Univalent vertices (dots). These have degree $+1$. \item Trivalent vertices, where all three adjoining edges have the same color. These have degree $-1$. \item
$2m$-valent vertices, where the adjoining edges alternate in color between two elements $s \ne t \in S$, and $m_{st}=m<\infty$. These have degree $0$. \end{itemize}

The \emph{degree} of a Soergel graph is the sum of the degree of each vertex, and the degree of each box. \end{defn}

\begin{figure} \label{thegenerators} \caption{The vertices in a Soergel graph} $\ig{1}{SoergGens}$ \end{figure}

When there is no ambiguity we refer to a Soergel graph merely as a ``graph," even though it is an isotopy class of embedded graph. We may also wish to discuss Soergel graphs on the planar disc, on an annulus, etc.

% \begin{defn} A \emph{double dot} is a boundary-less connected component of a graph consisting of two dots connected by a single edge. To each simple root $\a_s$ we associate the
% double dot colored $s$. To a monomial in the simple roots we may associate the corresponding disjoint union of double dots, itself a Soergel graph. In this fashion, we may associate
% to each element $f \in R$ a formal $\Bbbk$-linear combination of Soergel graphs with empty boundary. Here is an example where $f=\a_s^2 \a_t$.
% 
% \igc{1}{polyexample}
% 
% A \emph{Soergel graph with polynomials} is an isotopy class of Soergel graph which may also be decorated with boxes, each box labelled by a homogeneous element of $R$. We have already
% given an unambiguous way of viewing a Soergel graph with polynomials as a formal linear combination of Soergel graphs. However, we think of the boxes as being colorless, while each
% individual double dot has a color.

A \emph{boundary dot} is a connected component of a graph consisting of a dot connected to the boundary by a single edge. A Soergel graph that contains no dots or trivalent vertices is a standard graph, as in the previous chapter. However, we draw the edges as solid lines, not dashed lines.

\begin{defn} Let $\DC$ (or if there is ambiguity, $\DC_{S}$) denote the $\Bbbk$-linear monoidal category defined as follows. Objects are sequences $\un{w}$, sometimes denoted $B_{\un{w}}$, with monoidal structure given by concatenation. The space $\Hom_{\DC}(\un{w},\un{y})$ is the free $\Bbbk$-module generated by Soergel graphs with
bottom boundary $\un{w}$ and top boundary $\un{y}$, modulo the relations below. Hom spaces will be graded by the degree of the graphs, and all the relations below are homogeneous.

\emph{The polynomial relations:}
\begin{gather}
\label{alphais}
\begin{array}{c}
\tikz[scale=0.9]{
\draw[dashed,gray] (0,0) circle (1cm);
\draw[color=red] (0,0.5) -- (0,-0.5);
\node[draw,fill,circle,inner sep=0mm,minimum size=1mm,color=red] at (0,0.5) {};
\node[draw,fill,circle,inner sep=0mm,minimum size=1mm,color=red] at (0,-0.5) {};
} \end{array}
=
\begin{array}{c}
\tikz[scale=0.9]{
\draw[dashed,gray] (0,0) circle (1cm);
\node at (0,0) {$\alpha_s$};
} \end{array}, \\
%
\label{dotforcegeneral} 
\begin{array}{c}
\tikz[scale=0.9]{
\draw[dashed,gray] (0,0) circle (1cm);
\draw[color=red] (0,-1) to (0,1);
\node at (-0.5,0) {$f$};
} \end{array}
=
\begin{array}{c}
\tikz[scale=0.9]{
\draw[dashed,gray] (0,0) circle (1cm);
\draw[color=red] (0,-1) to (0,1);
\node at (0.5,0) {${\color{red}s}(f)$};
} \end{array}
+
\begin{array}{c}
\tikz[scale=0.9]{
\draw[dashed,gray] (0,0) circle (1cm);
\draw[color=red] (0,-1) to (0,-0.5);
\draw[color=red] (0,1) to (0,0.5);
\node[color=red,draw,fill,circle,inner sep=0mm,minimum size=1mm] at (0,-0.5) {};
\node[color=red,draw,fill,circle,inner sep=0mm,minimum size=1mm] at
(0,0.5) {};
\node at (0,0) {$\partial_{\color{red}s} f$};
} \end{array}.
\end{gather}

\emph{The one color relations:}
\begin{equation} \label{assoc1} \ig{1}{assoc1} \end{equation}	
\begin{equation} \label{unit} \ig{1}{unit} \end{equation}
\begin{equation} \label{needle} \ig{1}{needle} \end{equation}

\emph{The two color relations:} In order to simplify this presentation greatly, we will assume that our realization is balanced. For discussion of the general case, see Section \ref{nonbalanced}.

The color scheme depends slightly on the parity of $m=m_{st}<\infty$. We give one example of each relation for each
parity; the reader can guess the general form.

\begin{equation} \label{assoc2} \ig{1}{assoc2} \end{equation}
\begin{equation} \label{dot2m} \ig{1}{dot2m} \end{equation}

In equation \eqref{dot2m} above, the \emph{Jones-Wenzl morphism} $JW_{m-1}$ is a $\Bbbk$-linear combination of graphs constructed only out of dots and trivalent vertices. It will be
defined and discussed in the next section.

\emph{The three color relations:} It will be clear from the graphs which colors represent which indices. These relations are identical to those defined for $\DStd$ earlier, with the exception of
$H_3$.

For a triplet of colors forming a sub-Coxeter system of type $A_1 \times I_2(m)$ for $m<\infty$, we have

\begin{eqnarray} (A_1 \times I_2(m)) && \ig{1}{A1I2m} \label{A1I2m} \end{eqnarray}

A specific example, when $m=2$, is the case $A_1 \times A_1 \times A_1$:

\begin{eqnarray} (A_1 \times A_1 \times A_1) && \ig{1}{A1A1A1} \label{A1A1A1} \end{eqnarray}

The last three relations are for types $A_3$, $B_3$, and $H_3$ respectively, and are known as the \emph{Zamolodzhikov relations}. Unfortunately, $H_3$ is not complete.

\begin{eqnarray} (A_3) && \ig{1}{A3} \label{A3} \end{eqnarray}

\begin{eqnarray} (B_3) && \ig{1}{B3} \label{B3} \end{eqnarray}

\begin{eqnarray} (H_3) && 
\ig{1}{H3left}
- 
\ig{1}{H3right}
= 
\text{lower terms}
\label{H3} \end{eqnarray}

This concludes the definition of $\DC$. \end{defn}

\begin{remark} \label{incomplete} What it means to be a ``lower term" in \eqref{H3} was discussed in Section \ref{subsec-RexMoves}. Exactly what the lower terms are is currently unknown.
There is a correct answer, which is whatever makes the functor $\FC$ well-defined. We discuss how this could be addressed computationally below. \end{remark}

% We can not think of any possible applications for the precise formula anyway, unless you
% really really like type H. On the other hand, every diagram of these lower terms is assumed to be in light leaves form (see below). This does have certain implications for what can occur.
% ARE THERE ANY DOTS ON THE RIGHTMOST/LEFTMOST STRAND? DOES EVERY DIAGRAM FACTOR THROUGH $\le l(w_0) - 2$ strands??? \end{remark}

Note that (graded) Hom spaces are enriched in graded $R$-bimodules, since one can put a polynomial in a box and place it in the leftmost or rightmost region. We will see much later that all
hom spaces are free when considered as right or left $R$-modules, however this is far from clear at this stage. Because of \eqref{dotforcegeneral}, any diagram is equal to a linear combination of diagrams where polynomials only appear in the left-hand region. It is also an easy consequence of \eqref{dotforcegeneral} that, for $f \in R^s$, we have the \emph{polynomial sliding relation}:
\begin{equation} \label{dotslidegeneral} \ig{1}{dotslidegeneral} \end{equation}

Fix $k \ge 0$. Relations \eqref{unit} and \eqref{assoc1} imply that any two non-empty trees of a single color connecting $k$ points on the boundary are equal.

Let us assume Demazure Surjectivity, so that there exists some $\d$ with $\pa_s(\d)=1$. It is not difficult to use the one color relations to prove the decomposition \begin{equation}
\label{BsBsDecomp} B_s \ot B_s \cong B_s(1) \oplus B_s(-1). \end{equation} One simply needs to find two inclusion maps $i_1,i_2 \co B_s \to B_s \ot B_s$ of degree $-1,+1$ respectively, and
two projection maps $p_1,p_2$ of degree $+1,-1$ respectively, such that the usual relations are satisfied: $p_1 i_1 = p_2 i_2 = \1_{B_s}$, $p_1 i_2 = p_2 i_1 = 0$, and $\1_{B_s \ot B_s} =
i_1 p_1 + i_2 p_2$. One can choose these maps as follows. \igc{1}{iiidemp1}
This splitting is not canonical. In fact, the space of maps in degree $-1$ is one-dimensional, so that $i_1$ and $p_2$ are canonical. However, there may be many choices of $\d$.

\begin{remark} In fact, the choice of splitting is even more general than a choice of $\d$. For instance one may have
\begin{equation} \ig{1}{iiidemp2} \label{iiidemp2} \end{equation}
so long as $f + g = -\a_s$. Using \eqref{dotforcegeneral}, the example above has $f = \d - \a_s$ and $g= -\d$. However, this more general splitting exists even when $\hg$ does not satisfy Demazure Surjectivity: for instance, one could let $f = -\a_s$ and $g=0$. \end{remark}

\begin{remark} Using \eqref{dotforcegeneral}, one can replace two facing dots (a ``broken edge") with a sum of two diagrams having a complete edge, one with $\d$ on the left and the other
with $-s(\d)$ on the right. We call this procedure \emph{fusing} two dots. However, in the absence of Demazure Surjectivity, one can only fuse two dots up to a scalar, and the double
leaves ``basis" we define in the next chapter will no longer span all morphisms. Thus Demazure Surjectivity is still very important for proper behavior of the diagrammatic category. \end{remark}

See Remark \ref{twocolorrelncatfd} for the categorification of \eqref{rank2hecke}.

\begin{remark} \label{SCnotfinite} This remark is a diagrammatic analogue of Remark \ref{includingindexsets}. Given an inclusion of Coxeter systems $S_1 \subset S_2$, there is natural
functor $\iota \colon \DC_1 \to \DC_2$ which sends a Soergel graph for $S_1$ to itself, viewed as a Soergel graph for $S_2$ (reinterpreting boxes via the map $R_1 \to R_2$). As in the
bimodule case, this functor is not full for the trivial reason that $R_2$ is bigger than $R_1$, and thus $\DC_{S_2}$ has more boxes. However, it will follow from the main theorem of this
paper that $\iota$ is fully faithful after base change on the left from $R_{1}$ to $R_{2}$.

Suppose that a color $s$ does not appear on the boundary of a graph. The fact that $\iota$ is full after base change implies that we may manipulate the graph using our relations so that it
is in the span of graphs with polynomials where the color $s$ does not appear at all. This ``color removal" operation can be performed simultaneously for any number of colors. The proofs
used in type $A$ in \cite{EKh} actually provided a direct graphical algorithm for removing extraneous colors from a graph (in certain cases). Such an algorithm in general type would be
interesting, and remains an open problem. Nonetheless, we prove indirectly that extraneous colors can be removed by constructing a basis without them.

One can use these inclusions to define $\DC$ for Coxeter systems where $S$ is infinite, as a limit over the finite subsets of $S$. Any diagram contains finitely many colors and thus lies
inside a finite subcategory. The ring $R$ will no longer be Noetherian, but aside from that, Hom spaces will have all the nice finiteness properties (as free $R$-modules) that they have
when $S$ is finite. \end{remark}

\begin{remark} \label{generatorsinorder} As a monoidal category, $\DC$ can be defined with generators and relations. The generating morphisms are the vertices above, and cups and caps of
each color. However, relation \eqref{unit} implies that cups and caps can be constructed out of dots and trivalent vertices. Therefore, a full set of generators for $\DC$ is as listed in the introduction. 
% consists of:
% \begin{enumerate} \item For each color $s$, the dot oriented as a map $s \to \emptyset$. \item For each color $s$, the dot oriented as a map $\emptyset \to s$. \item For each color $s$,
% the trivalent vertex oriented as a map $ss \to s$. \item For each color $s$, the trivalent vertex oriented as a map $s \to ss$. \item For each pair of colors $s,t$ forming a finite
% dihedral group, the $2m$-valent vertex oriented as a map $\ubr{sts\ldots}{m_{st}} \to \ubr{tst\ldots}{m_{st}}$. \end{enumerate}
When constructing a functor to $\BSBim$, we will also need to check the unwritten isotopy relations. This is standard in diagrammatic categories; see, for instance, \cite{Lauda}. \end{remark}

%===========
\subsection{Jones-Wenzl morphisms} \label{subsec-JW}
%===========

Presenting the Bott-Samelson category in rank $2$ is the topic of \cite{EDihedral}. The calculations and proofs are too long to duplicate or fully discuss in this paper. If the reader is
willing to accept the Jones-Wenzl morphisms as black boxes\footnote{no pun intended!}, satisfying the properties stated in this section and the next, then the reader need not consult
\cite{EDihedral}. We assume the reader is familiar with quantum numbers and Temperley-Lieb algebras; background can be found in \cite{EDihedral}.

Fix $m<\infty$. Consider the Temperley-Lieb algebra $TL$ over $\ZM[z]$ with $m-1$ strands, where the circle is evaluates to $-z$. In any $\Bbbk[z]$-algebra we will use quantum
numbers to indicate the images of certain polynomials in $z$, where $[2]$ is the image of $z$. For instance, $[3]$ is the image of $z^2-1$. When we say that $q$ is specialized to a
primitive $2m$-th root of unity in a $\ZM[z]$-algebra, we mean that $m$ satisfies the minimal polynomial of $[2]$ at this root of unity, which equates to the statements that $[m]=0$, $[k]\ne 0$ for $0< k < m$, and $[m-1]=1$.

After inverting some quantum numbers, Temperley-Lieb algebras contain elements known as \emph{Jones-Wenzl projectors}, and the Jones-Wenzl projector on $m-1$ strands $JW_{m-1}$ is known to
be negligible and rotation-invariant when $q$ is specialized to a primitive $2m$-th root of unity. The Jones-Wenzl projector $JW_m$ can be defined whenever certain quantum binomial
coefficients are invertible. Using diagrammatics for the Temperley-Lieb algebra, the Jones-Wenzl projector can be described as a linear combination of crossingless matchings.

\begin{ex} $\ig{1}{JWinTL}$ \end{ex}

Any crossingless matching will divide the planar strip into $2(m-1)$ regions, which can be colored (say, red and blue) so that colors always alternate across a strand. This gives rise to
the \emph{two-colored Temperley-Lieb algebra}, which is an endomorphism algebra appearing inside a 2-category with 2 objects, red and blue. In this algebra, the value of a blue circle
surrounded by red and a red circle surrounded by blue can be different scalars; so we define this algebra over the ring $\ZM[x,y]$. Jones-Wenzl projectors still exist, but each comes in
two flavors with different coefficients, depending on whether blue or red appears on the far right. The other flavor is obtained by switching the colors and switching $x$ and $y$.

\begin{ex} $\ig{1}{JWin2TL}$ \end{ex}

Choose two indices $s,t \in S$. Given a 2-colored crossingless matching, we obtain a Soergel graph on the planar disc as follows: deformation retract each region into a tree composed out
of trivalent and univalent vertices, and color these trees appropriately. The result will always be a Soergel graph of degree $+2$. In order for this map to be well-defined, not just for
crossingless matchings but also for diagrams with embedded circles, one must specialize $\ZM[x,y]$ to $\Bbbk$ under the map sending $x \mapsto a_{st}$ and $y \mapsto a_{ts}$. The reader
should convince themselves that this makes sense, using \eqref{dotforcegeneral}, \eqref{alphais}, and \eqref{needle} until the scalar $\pa_s(\a_t) = a_{st}$ appears. Associated to the
Jones-Wenzl projector we have a linear combination of Soergel graphs, which we call the \emph{Jones-Wenzl morphism}. It comes in two color-flavors, as before.

\begin{ex} $\ig{1}{JWinSoergAlt}$ \end{ex}

The coefficients in these Jones-Wenzl projectors are rational functions in the 2-colored quantum numbers, as in Definition \ref{def:2qnum}. When $[m]_x=[m]_y=0$, one can show that
$JW_{m-1}$ is well-defined within $\ZM[x,y]$ (i.e. no denominators are necessary). If either $m$ is odd or if $m$ is even and $[m-1]=1$, $JW_{m-1}$ will be rotation-invariant (under
color-preserving rotations). If $m$ is even and $[m-1]=1$, it is also the case that rotating the left-blue-aligned Jones-Wenzl morphism by one strand yields the left-red-aligned
Jones-Wenzl morphism. If $m$ is odd, this holds if and only if $[m-1]_x = [m-1]_y = 1$; this is the feature behind the definition of a balanced realization. Using the examples above, the
reader should convince themselves of these rotational facts when $m=4$ (since $a_{st} a_{ts} = 2$) or when $m=3$ and $a_{st}=a_{ts}= \pm 1$. More discussion of the non-balanced case is
found in Section \ref{nonbalanced}.

The Jones-Wenzl Soergel graph above is not technically a morphism in $\DC$, because it lives on the planar disc, but it can be plugged into another diagram to produce a planar strip graph.
This is done in relation \eqref{dot2m}, where the rotation-invariant $JW_{m-1}$ is used. An implication of \eqref{dot2m} and \eqref{assoc2} is the following relation.
\begin{equation} \ig{1}{twocoloridemp} \label{twocoloridemp} \end{equation}
Note that here and henceforth, we will only use $JW_{m-1}$ when $m_{s,t}=m$, which we abbreviate as $JW$.

On the RHS of \eqref{twocoloridemp}, we took the Jones-Wenzl in the two-colored Temperley-Lieb algebra, transformed it into a degree $2$ Soergel graph on the disc, and then transformed it
again into a degree $0$ endomorphism of a color-alternating Bott-Samelson in $\DC$, by placing trivalent vertices on each side. This general procedure can be applied to any two-colored
crossingless matching, giving a map from the two-colored Temperley-Lieb algebra to a certain endomorphism ring in $\DC$. This map is a homomorphism. because $JW_{m-1}$ is idempotent in
$TL$, the RHS of \eqref{twocoloridemp} is also idempotent. The $2m$-valent vertex is thus ``half an idempotent."

The defining properties of the Jones-Wenzl projectors in $TL$ correspond to the following two properties in Soergel graphs. The first is that the coefficient of a certain graph is $1$, as
can be seen from the examples above. This is the graph which, when trivalent vertices are attached to make a degree $0$ morphism, becomes the identity map. The second property is called
\emph{death by pitchfork}:

\begin{equation} \ig{1}{JWpitch} \label{JWpitch} \end{equation}

Now let us note some properties of any $(m-1,m-1)$ crossingless matching, transformed into a Soergel graph. When viewed as a degree $+2$ graph on the disc, there will be a boundary dot on
at least one of any $m-1$ adjacent strands. When viewed as a degree $+0$ endomorphism of a length $m$ color-alternating object $B_s \ot B_t \ot \ldots$ in $\DC$, it can be given a
\emph{negative-positive decomposition}. That is, we say the initial \emph{width} of the map is $m$, the number of strands on the boundary. Reading from bottom to top, we may first apply
all negative maps (i.e. bottom boundary dots, merging trivalents, and caps) which are the maps which decrease the width. Then we can apply all positive maps (i.e. cups, splitting
trivalents, top boundary dots), increasing the width back up to $m$. The morphism will factor through a length $k$ color-alternating object $B_s \ot B_t \ot \ldots$ for some $k \le m$
which is the minimal width reached. Whenever the crossingless matching is not the identity map, we have $k < m$. These properties will be used in Chapter \ref{sec-LLLproof}.

\begin{remark} \label{twocolorrelncatfd} Consider the usual Temperley-Lieb category over a $\ZM[z]$-algebra where $[m]=0$ and $[k] \ne 0$ for $0 < k < m$. Then the Temperley-Lieb algebra
with $k$ strands is isomorphic to (an integral form of) the endomorphism ring of the $U_q(\mathfrak{sl}_2)$-representation $V_1^{\ot k}$ (at generic $q$), for any $0 \le k<m$. The direct
sum decomposition \[V_1^{\ot m-1} \cong V_d^{\oplus c_{m-1,d}} \] implies the existence of certain idempotents (and isomorphisms between their images within an isotypic component) which
produce this splitting. A similar statement can be made for the two-colored Temperley-Lieb category. We can transform these maps into degree $0$ morphisms in $\DC$, to obtain a direct sum
decomposition for any color-alternating $B_s \ot B_t \ot \ldots$ of length $\le m$. This can be used to categorify relation \eqref{rank2hecke}. See \cite{EDihedral} for more details.

However, we cannot assume in general that $[k] \ne 0$ for all $k < m$. For instance, this is false when we take a realization of a finite dihedral group, and view it as a non-faithful
realization of a larger dihedral group. Nonetheless, we will soon categorify the relation of Example \ref{rank2heckeEx}. \end{remark}

%===========
\subsection{The functor to bimodules} \label{subsec-functor}
%===========

We now fix a realization where Demazure Surjectivity holds. This ensures that $\Delta_s \in R \ot_{R^s} R$ is well-defined (see section \ref{subsec-polys}).

\begin{defn} Let $\FC$ be the $\Bbbk$-linear monoidal functor from $\DC$ to $\BSBim$ defined as follows. The object ${\un{x}}$ is sent to $B_{{\un{x}}}$. The images of the dots and trivalent vertices were given in the introduction, and correspond to the four structure maps of a Frobenius extension. In the introduction, one of the dots was sent to $1 \mapsto \frac{1}{2} (\a_s \ot 1 + 1 \ot \a_s)$, but this should be $1 \mapsto \Delta_s$ in general; whenever $\frac{1}{2}$ exists, these two expressions are equal. The bimodule image of the $2m$-valent vertex is given explicitly in \cite{EDihedral}, though in a convoluted and not particularly useful form. \end{defn}

\begin{claim} This definition gives a well-defined functor. \end{claim}

\begin{proof} Previous papers have done most of the work for this claim. The polynomial relations are obvious. Any other relation involving only a subset of colors can be checked in the
category $\DC_{S'}$, where $S'$ is the appropriate subset. Most of the relations (including the isotopy relations) involve at most 2 colors, and the dihedral case was checked in
\cite{EDihedral}. It remains to check the relations arising from rank 3 parabolics. The case of $A_3$ and $A_1 \times I_2(m)$ for $m=2,3$ was done already in \cite{EKh}, where general type
$A$ was completed. The check for $A_1 \times I_2(m)$ for other $m < \infty$ exactly parallels the proof for $m=2,3$ in \cite{EKh}, and is essentially trivial.

The case of $B_3$ was checked by computer. The case of $H_3$ can be checked by computer, once the appropriate relation is found. The computer check used a localization technique, which we
will discuss in the section after next. \end{proof}

%===========
\subsection{Localization}
%===========

\begin{defn} Let $\DC_Q$ denote the localization of $\DC$ at $Q$, which is to say that we allow boxes labelled by $f \in Q$ in the leftmost region, and require that they multiply as
in $Q$. Let $\FC_Q$ denote the functor $\DC_Q \to \BSBim_Q$ which extends $\FC$ under base change. \end{defn}

Because of arguments akin to Lemma \ref{localizationononeside}, this is the same as the category which allows boxes labelled by $f \in Q$ in any region, and allows
$f \in Q^s$ to slide across a line colored $s$.

The bimodule $B_s$ is equipped with short exact sequences \eqref{ses1} and \eqref{ses2} which split after localization to $Q$. Each short exact sequence ``splits" the other, in the sense
that following one map from \eqref{ses1} and then one map from \eqref{ses2} gives an endomorphism of $Q$ or of $Q_s$ which is multiplication by $\a_s$, and is therefore invertible. Two of
the maps we have already seen: they are the bottom and top boundary dot. We seek a calculus which mixes Soergel bimodules and standard bimodules, and which contains the other two maps. We
might draw them as follows.

\begin{equation} \ig{1}{bivalentgenerators} \label{bivalentgenerators} \end{equation}

\begin{defn} Let $\Kar(\DC_Q)$ (temporarily abusive notation) denote the \emph{non-cyclic(!)} biadjoint monoidal category generated on top of $\DC_Q$ as follows. Add new objects $Q_s$ for
each $s \in S$, called \emph{reflection indices}, whose identity morphisms we draw as dashed lines of the same color as $s$. Thus an object of $\Kar(\DC_Q)$ will have objects which are
sequences of normal indices and reflection indices. We typically disambiguate by writing $B_{\un{x},Q}$ for a sequence of normal indices and $Q_{\un{x}}$ for a sequence of reflection
indices. Allow an additional kind of vertex of degree $+1$, a bivalent vertex with one solid $s$ edge and one dashed $s$ edge. We also allow dashed cups and caps. Finally, impose these new
relations:

\begin{equation} \ig{1}{newidemprel1} \label{newidemprel1} \end{equation}
\begin{equation} \ig{1}{newidemprel2} \label{newidemprel2} \end{equation}
\begin{equation} \ig{1}{biadjointdashed} \label{biadjointdashed} \end{equation}
\begin{equation} \ig{1}{rotatebivalent} \label{rotatebivalent} \end{equation}
	
\end{defn}

The reader can check that these relations are consistent with \eqref{bivalentgenerators}. Note that we can define this category over $R$ instead of $Q$ if desired, or even in the absence
of Demazure Surjectivity. The remainder of this section will investigate $\Kar(\DC_Q)$ further, and prove the following theorem.

\begin{thm} Assume Demazure Surjectivity. This diagrammatic category is equivalent to the Karoubi envelope of $\DC_Q$. It is also equivalent to $\DStd_Q$, and to $\StdBim_Q$. \label{localdiagthm} \end{thm}

Unfortunately, $\Kar(\DC_Q)$ is not cyclic! It is worse than that the rotation of one bivalent vertex is not the other. The bivalent vertex itself is not invariant under 360 degree
twisting, being off by a sign. Any morphism with an even number of bivalent vertices will be cyclic (i.e. invariant under 360 degree twists), but this still does not imply that twisting by
180 degrees is equal to rotation by 180 degrees, because bivalent vertices do not behave this way consistently. An example can be seen below in \eqref{twist2mwbivalent}. At least different
representatives of the same isotopy class only differ by a sign.

\begin{remark} \label{geometrynocyclic} If one is familiar with the geometric underpinning of Soergel bimodules, then it should not be terribly off-putting that these maps are not cyclic.
One expects cyclicity whenever one analyzes convolution between perverse sheaves, because the procedure of taking the biadjoint of a sheaf is a functor. Soergel bimodules are the
equivariant (derived) global sections of the semisimple perverse sheaves on the flag variety, and taking global sections is a well-behaved functor on this semisimple category, so that
cyclicity happens to be preserved. However, taking global sections of non-semisimple sheaves tends to forget structure, and break the compatibility with the biadjunction functor. If one
more appropriately models non-semisimple perverse sheaves as complexes of semisimple sheaves, then these complexes will be biadjoint in a cyclic way.

For instance, the standard bimodule $R_s$ corresponds to two different, mutually biadjoint perverse sheaves: the shriek and the star extension of the constant sheaf on $\mathbb{P}^1$ minus
a point. These two different perverse sheaves have two different resolutions in terms of semisimple perverse sheaves, and these descend to \eqref{ses1} and \eqref{ses2}. The two complexes
of Soergel bimodules which are quasi-isomorphic to $R_s$ are biadjoint in a cyclic way. The bimodule $R_s$ itself is self-biadjoint. However, the compatibility between these two
biadjunction structures is broken. \end{remark}

We call an isotopy class of graphs as above \emph{mixed graphs}. A mixed graph only represents a morphism in $\Kar(\DC_Q)$ up to sign. A mixed graph has some solid edges and some dashed
edges. A mixed graph without any dashed edges is a Soergel graph or \emph{solid graph}, and a mixed graph without any solid edges is a standard graph or \emph{dashed graph} (we will expand
what counts as a dashed graph soon). (Isotopy classes of) solid graphs and dashed graphs do unambiguously represent a morphism, without any sign issues, because they have no bivalent
vertices.

Now let us do some calculations. Using \eqref{newidemprel1} and \eqref{newidemprel2}, one can easily produce the following equalities after multiplication by $\a_s$. Since we're working
over $Q$, we may divide by $\a_s$.

\begin{equation} \ig{1}{dotbi} \label{dotbi} \end{equation}

\begin{remark} If defining this diagrammatic category over $R$ instead of $Q$, one should add \eqref{dotbi} to the list of relations. \end{remark}

The new bivalent vertices give rise to an idempotent $\ig{1}{dashedidemp}$, which is complementary to the idempotent $\ig{1}{brokenidemp}$ which already existed in $\DC_Q$. Therefore,
$B_{s,Q} \cong Q_s(1) \oplus Q(1)$. (We include the gradings for those who wish to use the homogeneous fraction field rather than the full fraction field. Remember that $Q \cong Q(2)$ so that the grading lives in $\Z/2\Z$.)  By convention we tend to include $\frac{1}{\a_s}$ in the projection map, rather than the inclusion map.

\begin{equation} \ig{1}{localdecomp} \label{localdecomp} \end{equation}
We can take a line and either ``break" it or ``dash" it. How very violent!

\begin{claim} The reflection indices behave like standard modules with respect to polynomials. In other words, we have \eqref{slidepolythrustd}. \end{claim}

\begin{proof} Place a polynomial $f$ on the left side of the diagrams in \eqref{newidemprel1}. Use the polynomial forcing rules to force $f$ through the solid line in the middle. Any term
where the line breaks is zero by \eqref{dotbi}. The remaining term has $s(f)$ on the right hand side instead. Dividing by $\a_s$ gives the desired equality. \end{proof}

\begin{claim} The dashed cups and caps are redundant, being equal to the following maps. The only relations necessary are \eqref{newidemprel1} and \eqref{newidemprel2}.
\begin{equation} \ig{1}{dashedcupscaps} \label{dashedcupscaps} \end{equation} \end{claim}

\begin{proof} Equation \eqref{dashedcupscaps} follows from \eqref{newidemprel1} by adding a cap or cup and using \eqref{rotatebivalent}. Conversely, if we only use \eqref{dashedcupscaps}
as a definition of the dashed cups and caps, it is a simple calculation to check \eqref{rotatebivalent} and \eqref{biadjointdashed}. We give a sample computation here.
\igc{1}{cupcalcexample} \qedhere
\end{proof}

Fix a category $\CC$ specified with generators and relations, an object $M \in \CC$, and an idempotent $e \in \End(M)$. Let $\CC(M,e)$ denote the partial idempotent completion which
formally adds the image of $e$ as a new object. Let us call this new object $X$. It is easy to give a presentation of $\CC(M,e)$ by generators and relations. One adds a new object $X$ and
two new generators, a map $\iota \colon X \to M$ and a map $\pi \colon M \to X$. One adds two new relations, which state that $\iota \pi = e \in \End(M)$ and $\pi \iota = \1 \in \End(X)$.
This presentation clearly gives a category $\CC(M,e)$ equipped with a fully faithful map $\CC \to \CC(M,e)$ which has all the desired properties. If the image of $e$ is already an object
in $\CC$ then $X$ will be isomorphic to it. Similarly, if we wish to adjoin a set of new summands $\{X_\b\}$, we need only add inclusions $\{\iota_\b\}$ and projections $\{\pi_\b\}$ with
the relations above for each $\b$. Similar statements can be made about monoidal categories with monoidal presentations, graded categories etc.

Now it is clear that $\Kar(\DC_Q)$ is the category obtained from $\DC_Q$ by adjoining the complements of $Q$ in $B_{s,Q}$, for each $s$. The bivalent vertices are the new maps $\iota$
and $\pi$, and the two relations on $\iota$ and $\pi$ correspond to \eqref{newidemprel1} and \eqref{newidemprel2}. To show that $\Kar(\DC_Q)$ is the Karoubi envelope of $\DC_Q$, it is
enough to show that $\Kar(\DC_Q)$ is idempotent complete. Because each $B_{s,Q}$ decomposes into $Q_s$ and the monoidal identity, it is clear that any object in $\DC_Q$ is isomorphic to a
direct sum of sequences consisting solely of reflection indices. Thus we need only show that any sequence consisting of reflection indices represents an indecomposable object. This will be
implied once we show that this diagrammatic category is equivalent to $\DStd_Q$.

We do some more computations in preparation for the proof of this equivalence. Suppose that one takes a $2m$-valent vertex and places $m$ consecutive bivalent vertices on it (say,
on bottom).

\igc{1}{2mwbivalent}

If one places a dot on top of this diagram, we get $0$. This is because we can use \eqref{dot2m} to replace the $2m$-valent vertex with the Jones-Wenzl morphism, and the Jones-Wenzl
morphism must have a dot on one of the bottom exits. This dot will then hit the bivalent vertex, yielding zero by \eqref{dotbi}. However, each line on top decomposes into two idempotents
as in \eqref{localdecomp}, and only the dashed idempotent survives. Thus if there are ever $m$ consecutive bivalent vertices on a $2m$-valent vertex, we may as well assume that all $2m$
are present.

If we place $2m$ bivalent vertices around a $2m$-valent vertex, we get a morphism of degree $2m$. In order to get a morphism of degree $0$ we should divide by a polynomial. In other words,
half the bivalent vertices should be inclusions and half projections, and one half should be paired with $\frac{1}{\a_s}$. So consider the following diagram with $m$ strands.

\igc{1}{w02mdashed}

This is a reduced expression for $w_0$, the longest element of $W_{s,t}$. Using \eqref{newidemprel1}, we get a sequence of vertical dashed lines, with $\a_s$ to the left of each line
colored $s$. If we pull all these polynomials to the far left region using \eqref{slidepolythrustd}, we get $\rho_{s,t}$, the product of all $m$ positive roots corresponding to reflections
in $\langle s, t \rangle$ (for a definition of positive roots in a dihedral group, see \cite{EDihedral}). Note that $w_0$ sends the set of positive roots to the set of negative roots, so
that $w_0(\rho)=(-1)^m \rho$.

A simple calculation using \eqref{rotatebivalent} shows that

\begin{equation} \ig{1}{twist2mwbivalent} \label{twist2mwbivalent} \end{equation}

Therefore

\begin{equation} \ig{1}{2mstdiscyclic} \label{2mstdiscyclic} \end{equation}

Thus the RHS of \eqref{2mstdiscyclic} is a degree $0$ map which \emph{is} cyclic, and is a perfect candidate for the $2m$-valent vertex in $\DStd_Q$.

\begin{defn} Let $\StdF$ denote the functor from $\DStd_Q$ to $\Kar(\DC_Q)$ defined herein, called the \emph{standardization functor}. On objects, it sends a sequence of indices to the
corresponding sequence of reflection indices. On morphisms, it sends caps and cups to dashed caps and cups, and it sends the $2m$-valent vertex to the morphism in \eqref{2mstdiscyclic}.
\end{defn}

We draw the image of the $2m$-valent vertex in $\DStd_Q$ as a dashed $2m$-valent vertex. From the above, it should be easy to check that

\begin{equation} \ig{1}{slidebivalentthru2m} \label{slidebivalentthru2m} \end{equation}

\begin{prop} The functor $\StdF$ is well-defined, and is an equivalence of categories. \label{StdFanequiv} \end{prop}

This proposition implies all of Theorem \ref{localdiagthm} except for the connection to $\StdBim$. We have already showed that $\StdF$ is essentially surjective, because $\Kar(\DC_Q)$ is
additively generated by reflection indices. We need to check that $\StdF$ is well-defined, full and faithful.

\begin{proof}[Proof that $\StdF$ is well-defined.] We have already checked the isotopy relations in $\StdBim$, because of \eqref{biadjointdashed} and \eqref{2mstdiscyclic}. We have also checked
polynomial-sliding. Relation \eqref{capcupstd} follows as below (we used $\d = \frac{\a_s}{2}$, but any $\d$ will work).

\igc{1}{capcupstdproof}

The proof of \eqref{circlestd} is easy, and we leave it as an exercise.

The proofs of \eqref{2misomstd} and the three color relations all follow from the same method. Take relations \eqref{twocoloridemp} and the three color relations in $\DC$, and place a
bivalent vertex below every strand on bottom. Doing so will kill any diagram with a bottom boundary dot, including all the non-identity diagrams in the Jones-Wenzl projector in
\eqref{twocoloridemp}, and all the lower terms in the $H_3$ relation \eqref{H3}. We ignore all those terms. For any diagram composed entirely out of $2m$-valent vertices for various $m$,
we can use \eqref{slidebivalentthru2m} to pull the ``dashed-ness" from bottom to top, until the entire diagram is dashed except with bivalent vertices at the top. The result is precisely
the corresponding relation in $\DStd$, with bivalent vertices on top. Bivalent vertices are invertible, so this checks the relation. \end{proof}

Now we can apply any relations in $\DStd$ to dashed diagrams in $\Kar(\DC_Q)$.

We see that surrounding a $2m$-valent vertex with bivalent vertices yields (up to polynomial) the dashed version of the map. However, there can be no dashed version of the dot or the
trivalent vertex, since there are no maps between standard bimodules when they do not express the same element of $W$. Unsurprisingly, surrounding a dot or a trivalent vertex with bivalent
vertices is zero. For the dot this is \eqref{dotbi}. It is not too hard to show that

\igc{1}{trivalentlemma}

Therefore

\begin{equation} \ig{1}{dashedtriiszero} \label{dashedtriiszero} \end{equation}

\begin{lemma} \label{reduceinKarDCQ} Every morphism in $\Kar(\DC_Q)$ is a linear combination of graphs for which: \begin{itemize} \item The only dots appearing are boundary dots. \item There is a single polynomial in the leftmost region, and no other connected components without boundary. \item There are no trivalent
vertices. \item There are no solid $2m$-valent vertices, only dashed $2m$-valent vertices. \end{itemize} \end{lemma}

\begin{proof} Consider a diagram in $\Kar(\DC_Q)$. One can use \eqref{localdecomp} to either break or dash any edge. The reduction goes as follows: \begin{enumerate} \item Counting regions
separated by a dashed line as a single \emph{area}, we can assume there is a single area. This can be done by taking a solid line which separates two areas and either breaking it or
dashing it. \item We can assume that every $2m$-valent vertex appears only in its dashed version. If there are any solid $2m$-valent vertices, break or dash one of its edges. If broken,
one can use \eqref{dot2m} to remove the $2m$-valent vertex. If dashed, now break or dash the next edge. Eventually we can assume every edge is dashed, any thus can replace the $2m$-valent
vertex with its dashed version (up to multiplication by a polynomial in some region). \item We can assume there are no trivalent vertices. If there are any trivalent vertices, break or
dash one of its edges. If broken, one can use \eqref{unit} to remove the trivalent vertex. If dashed, now break or dash the next edge. One of the three edges must be broken, or the result
is zero by \eqref{dashedtriiszero}. \item We can assume there are no dots except for boundary dots. Any dot not connected to the boundary must be connected to another dot (and so becomes a
box), or to bivalent vertex (and so becomes zero). \item Any remaining solid line can only run into the boundary or into a bivalent vertex (it can not form a circle, for this would create
a second area). Any solid line meeting two bivalent vertices can be removed by \eqref{newidemprel1}, yielding a continuous dashed line (up to a polynomial). Thus any connected component
without boundary is either a box or is a purely dashed diagram. \item We can assume all boxes are in the leftmost region. This is because there is a single area, and we can slide
polynomials through dashed lines using \eqref{slidepolythrustd}. Thus the polynomials merge into a single box. \item Any remaining closed component can be removed. This uses the fact that
closed diagrams in $\DStd$ reduce to the empty diagram. \end{enumerate} \end{proof}

\begin{cor} $\StdF$ is full. \end{cor}

\begin{proof} Apply the reduction of the lemma to a map between reflection indices. There can be no boundary dots. A solid edge can not connect to anything except a
bivalent vertex any more. Thus every solid edge can be removed with \eqref{newidemprel1} (at the cost of adding a polynomial). Any map with only dashed edges is clearly in the image of
$\StdF$. \end{proof}

We can define a functor $\Kar(\DC_Q) \to \StdBim$ extending $\FC$. This functor acts on bivalent vertices as in \eqref{bivalentgenerators}. Clearly this functor intertwines $\StdF$ and
$\FC_\std$. Since $\FC_\std$ is an equivalence, this would imply that $\StdF$ is faithful, and is therefore an equivalence. This concludes the proof of Proposition \ref{StdFanequiv} and Theorem \ref{localdiagthm}

In the next section we construct a quasi-inverse for $\StdF$, giving a proof that $\StdF$ is faithful without needing the functor $\FC$ to bimodules.

%===========
\subsection{Computation using localization}
\label{subsec-compute}
%===========

Each object $B_{\un{x}}$ in $\DC$ splits in $\Kar(\DC_Q)$ as the sum of $Q_\eb$ over all subsequences $\eb$ of ${\un{x}}$. Here, $Q_\eb$ denotes the object which is the tensor product of $Q_{{\un{x}}_k}$
when $\eb_k=1$ and $\1$ when $\eb_k=0$. Given a graph expressing a morphism from a sequence $B_{\un{x}}$ of length $d$ to a sequence $B_{{\un{x}}'}$ of length $d'$, we can localize to obtain a $2^{d'} \times 2^d$ matrix of maps
between reflection sequences. Since Hom spaces between standard bimodules are always either rank 1 or 0, this matrix is actually populated with polynomials in $Q$, and is fairly sparse
(because many Hom spaces are zero a priori). Computing any term in this matrix consists of applying the appropriate projection and inclusion maps to the top and bottom of the graph, and
using the diagrammatics of $\Kar(\DC_Q)$ to reduce the graph to a dashed graph with the desired polynomial on the left. That is, given $({\un{x}},\eb)$ and $({\un{x}}',\eb')$, the coefficient of the
map $Q_\eb \to Q_{\eb'}$ of a map $\phi$ is given by reducing the following diagram $\phi_\eb^{\eb'}$.

\igc{2}{localizingamap}

There is no sign issue in this convention, even though a mixed graph represents a morphism only up to sign. For any Soergel graph $\phi$ we start with, we can choose a representative of
the isotopy class, and add idempotents as above to obtain a specific representative of the mixed graph. Different representatives of $\phi$ will give different mixed graphs, but they
differ only by isotopy of the solid part of the graph, and therefore have the same sign when viewed as morphisms in $\Kar(\DC_Q)$.

Note that by \eqref{slidebivalentthru2m}, the coefficient associated to $\eb=(1,1,\ldots,1)$ and $\eb' = (1,1,\ldots,1)$ of a $2m$-valent vertex is precisely $1$. Therefore, for any rex
move, this ``highest" coefficient will be $1$.

One can check if two maps in $\DC$ are equal by computing these two matrices and comparing the terms. More combinatorially, one can compute once and for all the ``basic" matrices
attached to the generating morphisms. For instance, the dot gives a $2 \times 1$ matrix, the trivalent vertex a $4 \times 2$ matrix, and the $2m$-valent vertex a $2^m \times 2^m$ matrix.
Computing a more general map consists merely of multiplying these basic matrices and annoying bookkeeping. Computers excel at such tasks.

We have not yet proven it, but the passage $\DC \to \DC_Q$ is faithful. Not knowing this, the equality of two matrices only implies that the original maps are equal modulo $R$-torsion.
However, location is injective in the bimodule world, because Hom spaces are free left $R$-modules. In other words, all $R$-torsion is in the kernel of the functor $\FC$. Therefore we can
calculate whether two Soergel graphs have the same image under $\FC$ by localizing $\DC \to \Kar(\DC_Q)$ and computing the matrices above. This is a powerful tool.

For instance, we want to know whether both sides of the $B_3$ relation \eqref{B3} correspond to the same map between Bott-Samelson bimodules. One need only compute by hand the matrices
associated to the 4-valent, 6-valent, and 8-valent vertices, and then plug in two appropriate formulae into a computer. Checking that two sparse matrices of size $2^8 \times 2^8$ are equal
is trivial. If we knew what the $H_3$ relation \eqref{H3} should be, then checking it would require calculating the 10-valent vertex (which is easy), and then computation of a pair of
$2^{15} \times 2^{15}$ matrices (which is quite time-consuming). Unfortunately, backsolving for the coefficients in the H3 relation would require doing linear algebra with a large number
$N$ of unknowns (at least $70$). There is an equation for each nonzero term in a sparse $2^{15} \times 2^{15}$ matrix, and the coefficients come from $N$ different such matrices which need
to be computed. Neither the author's computers nor their brains appear to be up to the task.

Finally, we provide the alternative proof of faithfulness.

\begin{prop} $\StdF$ is faithful. \label{stdfisfaithful2}\end{prop}

\begin{proof} In fact, the techniques we have developed allow us to construct a quasi-inverse for $\StdF$. We construct a functor $\GC$ from $\Kar(\DC_Q)$ to $\DStd_Q$ as follows. Let
$\GC$ send reflection sequences to themselves, and normal sequences to the corresponding formal direct sum of reflection sequences. On morphisms, $\GC$ will send standard graphs to
themselves. The generating morphisms which are not standard graphs are sent as follows: a bivalent vertex is sent to the appropriate inclusion or projection; a dot or trivalent vertex or
$2m$-valent vertex is sent to the appropriate matrix of standard diagrams. It is trivial to check that all the relations hold, so this functor is well-defined, and is obviously a
quasi-inverse to $\StdF$. \end{proof}

Remember that diagrams and their linear combinations are only useful for displaying morphisms between tensor products of generators, not for \emph{direct sums} of those. To talk about a
morphism between direct sums, we need to use matrices of diagrams. Thus we do not expect there to be an actual diagram in $\DStd_Q$ corresponding to a solid graph.

This proof that $\Kar(\DC_Q) \cong \DStd_Q$ was entirely diagrammatic, and no mention of bimodules was required.


%======================
\subsection{Unbalanced realizations}
\label{nonbalanced}
%======================

Now we discuss the diagrammatic alterations which must be made to accommodate the case of non-balanced realizations. This was discussed for dihedral groups in the appendix to
\cite{EDihedral}, where two separate diagrammatic conventions are proposed to deal with the new bookkeeping required. We follow the second suggested convention from that appendix.

Fix a dihedral parabolic subgroup with $m=m_{st} < \infty$. Relation \eqref{dot2m} implies that there is a close connection between $2m$-valent vertices and two-colored Jones-Wenzl
projectors. However, this begs the question: which Jones-Wenzl projector? After all, the left-blue-aligned Jones-Wenzl projector is not equal to the rotation of the left-red-aligned one. A
careful examination of \eqref{dot2m} and \eqref{twocoloridemp} shows that different choices of Jones-Wenzl projector must be made, depending on the location of the dot and the orientation
of the $2m$-valent vertex.

If some rescaling of \eqref{dot2m} is to hold, for any positioning of the dot and the $2m$-valent vertex, then the Jones-Wenzl projector must satisfy the death by pitchfork property
\eqref{JWpitch}. In particular, $JW_{m-1}$ must be negligible, and must have some rotational eigenvalue. This is not possible when the realization is even-unbalanced (see \cite{EDihedral}
for more details). While one may be able to design a diagrammatic calculus for even-unbalanced realizations, we will not attempt to do so. When $m$ is even and $[m-1]=1$, all the two-color
relations above hold as stated.

The $2m$-valent vertex is supposed to correspond, under the functor $\FC$, to some non-zero morphism between bimodules, living in a one-dimensional space of morphisms. Let $\un{w}_{(s)}$
denote the reduced expression $\ldots tsts$ of length $m$ ending in $s$, and let $\un{w}_{(t)}$ denote the reduced expression $\ldots stst$ of length $m$ ending in $t$. There is a unique
bimodule map $f_{(s)} \co BS(\un{w}_{(t)}) \to BS(\un{w}_{(s)})$ which sends the 1-tensor to the 1-tensor, and a unique map $f_{(t)} \co BS(\un{w}_{(s)}) \to BS(\un{w}_{(t)})$ which does
the same. However, when the realization is not balanced, these maps are not rotations of each other (by one strand), and this is the underlying issue. However, these maps are individually
invariant under color-preserving rotations (just like $JW_{m-1}$, even in the odd-unbalanced case), so we may draw each unambiguously as some kind of $2m$-valent vertex. We label the
vertices $(s)$ or $(t)$ to distinguish the two. (We only draw the case when $m$ is odd, but nothing prevents drawing the even case too.)

\igc{1}{funky2m}

The purple strand is meant to encode the appropriate sequence of alternating red and blue strands. To reiterate, when the blue-centered $2m$-valent vertex is oriented such that its
upper-right strand is blue (as in the picture above), then it corresponds under $\FC$ to a morphism which preserves the 1-tensor. When the blue-centered $2m$-valent vertex is oriented such
that its upper-right strand is red, one differs from this map by an invertible scalar $\l$. In \cite{EDihedral}, this scalar $\l = [m-1]_y$ is discussed at additional length. Therefore,
one has the rotational relation: \begin{equation} {
\labellist
\small\hair 2pt
 \pinlabel {$\l$} [ ] at 65 24
 \pinlabel {$\l^{-1}$} [ ] at 194 23
\endlabellist
\centering
\ig{1}{funky2mrotate}
} \end{equation}

When we place a dot on a $2m$-valent vertex, we obtain a relation similar to \eqref{dot2m}. \begin{equation} {
\labellist
\small\hair 2pt
 \pinlabel {$JW_{m-1}$} [ ] at 92 93
 \pinlabel {$JW_{m-1}$} [ ] at 236 93
 \pinlabel {$JW_{m-1}$} [ ] at 92 20
 \pinlabel {$JW_{m-1}$} [ ] at 236 20
\endlabellist
\centering
\ig{1}{dot2malternate}
} \end{equation} In each case, the version of the
Jones-Wenzl projector used is the one whose coefficient of the identity is $1$, when color-aligned as above. By rotating these pictures, one knows how to place a dot on any strand in
either $2m$-valent vertex. To check that these relations make sense, observe that both sides send the 1-tensor to the 1-tensor.

To give the correct versions of \eqref{assoc2} and \eqref{twocoloridemp}, one should label the $2m$-valent vertices in such a way that both sides preserve the 1-tensor. \begin{equation}
\ig{1}{twocolorassocalternate} \end{equation} \begin{equation} \ig{1}{twocoloridempalternate} \end{equation} There are additional, color-switched versions of each of these relations.

Finally, to give the correct version of the three-color relations, one should again ensure that both sides preserve the 1-tensor. Here is the $A_3$ relation. \begin{equation}
\ig{1}{A3altalt} \end{equation}

Note that, when working with these diagrams, there is a scalar ambiguity that appears when defining a rex move, determined by the choice of central color on each $2m$-valent vertex. Our
convention is that one will always choose a coloring so that the 1-tensor is preserved by the rex move.

For the remainder of this paper, we will work solely with the balanced case, and thus will not need the extra confusion of labeled $2m$-valent vertices. In the rest of the paper, it is
roughly the case that diagrams are only important ``up to scalar." That is, we will be asking whether certain morphisms span or are linearly independent, and these concepts are not
affected by rescaling. It should not be hard to convince oneself that the remainder of this paper applies equally well to the odd-unbalanced case.